% ============================================================
% Apéndice B — Resultados complementarios
% ============================================================

\chapter{Resultados complementarios}
\label{ap:resultados}

Este apéndice recoge las tablas completas de los experimentos del
Capítulo~\ref{cap:resultados}, con todas las pérdidas, post-procesos y
medidas. Todas las cifras son de val~$B$ salvo que se indique lo
contrario.

% ------------------------------------------------------------
\section{Post-proceso clásico}
\label{sec:anx-post}
% ------------------------------------------------------------

La Tabla~\ref{tab:anx-post} compara la búsqueda en rejilla y con Optuna
de los dos post-procesos clásicos sobre la baseline (semilla 42). Ambas
búsquedas llegan a valores parecidos, lo que indica que la rejilla ya
cubría la zona de buenos valores; en el trabajo se usa Optuna. Al subir
el umbral sin filtro de área, los falsos positivos aumentan
ligeramente, posiblemente porque algunas componentes se dividen.

\begin{table}[htbp]
\centering
\small
\begin{tabular}{@{}lccccccc@{}}
\toprule
Post-proceso & $t$ & $a$ & PQ & SQ & RQ & TP & FP / FN\\
\midrule
Sin post-proceso & 0,50 & 0 & 0,311 & 0,638 & 0,488 & 495 & 665 / 374\\
B1 (rejilla) & 0,65 & 0 & 0,314 & 0,647 & 0,485 & 494 & 674 / 375\\
B1 (Optuna) & 0,66 & 0 & 0,313 & 0,646 & 0,484 & 494 & 678 / 375\\
B2 (rejilla) & 0,65 & 25 & \textbf{0,375} & 0,649 & 0,577 & 473 & 297 / 396\\
B2 (Optuna) & 0,66 & 29 & 0,373 & 0,650 & 0,574 & 458 & 269 / 411\\
\bottomrule
\end{tabular}
\caption{Post-proceso de la baseline (semilla 42) en val~$B$. $t$ es el
umbral y $a$ el área mínima en píxeles a $512\times512$, ambos ajustados
en val~$A$.}
\label{tab:anx-post}
\end{table}

% ------------------------------------------------------------
\section{Comparación de pérdidas}
\label{sec:anx-perdidas}
% ------------------------------------------------------------

La Tabla~\ref{tab:res-perdidas} da el PQ de las diez pérdidas, en val~$A$
y val~$B$, sin post-proceso y con B2. La combinación de BCE y clDice sin
Dice colapsa. El mejor modelo en $A$ es el de la primera época (la sexta
con Topograph) y su PQ no llega a $0{,}05$ ni con post-proceso. Las
Tablas~\ref{tab:res-asignaciones} y~\ref{tab:res-fiabilidad} dan los
tipos de error y la fiabilidad de las detecciones. De ellas salen las
Figuras~\ref{fig:prec-exh} y~\ref{fig:tipos-error}.

Dos observaciones complementan las del capítulo. En partidos y unidos,
clDice es tan eficaz como Topograph, lo que concuerda con que premia la
continuidad de las estructuras alargadas, mientras que IoU, sin sesgo
topológico, llega a 123 filamentos partidos. Con B2, los filamentos
sin predicción suben en todos los modelos, hasta 201--301, y
IoU+Topograph es, tras Topograph (s43), el modelo con menos predicciones
sin filamento (54).

\begin{table}[htbp]
\centering
\small
\begin{tabular}{@{}lccccc@{}}
\toprule
 & & \multicolumn{2}{c}{Sin post-proceso} & \multicolumn{2}{c}{B2 (Optuna)}\\
\cmidrule(lr){3-4}\cmidrule(lr){5-6}
Pérdida & Época & PQ $A$ & PQ $B$ & PQ $A$ & PQ $B$\\
\midrule
Baseline (s42) & 8 & 0,297 & 0,311 & \textbf{0,364} & \textbf{0,373}\\
Baseline (s43) & 8 & 0,276 & 0,287 & 0,348 & 0,357\\
IoU & 15 & 0,279 & 0,285 & 0,351 & 0,351\\
clDice & 13 & 0,307 & 0,313 & 0,351 & 0,354\\
IoU+clDice & 14 & 0,308 & 0,322 & 0,345 & 0,363\\
clDice sin Dice & 1 & 0,012 & 0,009 & 0,062 & 0,036\\
\midrule
Topograph (s42) & 17 & \textbf{0,322} & \textbf{0,327} & 0,354 & 0,354\\
Topograph (s43) & 12 & 0,298 & 0,322 & 0,336 & 0,355\\
IoU+Topo & 16 & 0,315 & 0,315 & 0,346 & 0,346\\
clDice+Topo & 18 & 0,307 & 0,317 & 0,344 & 0,348\\
IoU+clDice+Topo & 18 & 0,297 & 0,318 & 0,333 & 0,336\\
clDice sin Dice+Topo & 6 & 0,025 & 0,014 & 0,080 & 0,045\\
\bottomrule
\end{tabular}
\caption{PQ de cada pérdida en val~$A$ y val~$B$, sin post-proceso
(umbral $0{,}5$, sin filtro de área) y con B2 ajustado en val~$A$ para
cada ejecución. La «Época» es la elegida en val~$A$ y la semilla se fija en 42 salvo que
se indique.}
\label{tab:res-perdidas}
\end{table}

\begin{table}[htbp]
\centering
\small
\begin{tabular}{@{}lcccccc@{}}
\toprule
Pérdida & Partidos & Unidos & P.\,+\,U. & Sin pred. & Sin fil. & S.\,P.\,+\,S.\,F.\\
\midrule
\multicolumn{7}{l}{\emph{Sin post-proceso}}\\
Baseline (s42) & 92 & \textbf{12} & 104 & 140 & 319 & 459\\
Baseline (s43) & 111 & 14 & 125 & \textbf{100} & 482 & 582\\
IoU & 123 & 13 & 136 & 119 & 380 & 499\\
clDice & 73 & 17 & 90 & 122 & 278 & 400\\
IoU+clDice & \textbf{57} & 15 & \textbf{72} & 124 & 261 & 385\\
Topograph (s42) & 78 & 13 & 91 & 183 & 169 & 352\\
Topograph (s43) & 68 & 13 & 81 & 224 & \textbf{110} & \textbf{334}\\
IoU+Topo & 79 & 13 & 92 & 199 & 142 & 341\\
clDice+Topo & 65 & 15 & 80 & 180 & 159 & 339\\
IoU+clDice+Topo & 63 & 18 & 81 & 197 & 161 & 358\\
\midrule
\multicolumn{7}{l}{\emph{B2 (Optuna)}}\\
Baseline (s42) & 43 & \textbf{12} & 55 & 250 & 68 & 318\\
Baseline (s43) & 37 & 13 & 50 & 275 & 88 & 363\\
IoU & 45 & 13 & 58 & 247 & 95 & 342\\
clDice & 39 & 15 & 54 & 213 & 119 & 332\\
IoU+clDice & 34 & 15 & 49 & \textbf{201} & 110 & \textbf{311}\\
Topograph (s42) & 41 & \textbf{12} & 53 & 217 & 108 & 325\\
Topograph (s43) & \textbf{26} & 14 & \textbf{40} & 300 & \textbf{49} & 349\\
IoU+Topo & 36 & 13 & 49 & 301 & 54 & 355\\
clDice+Topo & 37 & 14 & 51 & 240 & 82 & 322\\
IoU+clDice+Topo & 43 & 15 & 58 & 248 & 82 & 330\\
\bottomrule
\end{tabular}
\caption{Tipos de error en val~$B$, sumados sobre todos los pares
anotador--imagen. \emph{Partidos}: filamentos de referencia tocados por
dos o más predicciones. \emph{Unidos}: predicciones que tocan dos o más
filamentos. \emph{Sin pred.}: filamentos que ninguna predicción toca.
\emph{Sin fil.}: predicciones que no tocan ningún filamento. «Tocar»
significa solaparse en al menos un píxel ($\mathrm{IoU}>0$), como en la
definición oficial. P.\,+\,U.\ y S.\,P.\,+\,S.\,F.\ son las sumas que
miden los errores de estructura y de presencia
(Sección~\ref{sec:evaluacion}). En negrita, el menor valor de cada
columna dentro de cada bloque.}
\label{tab:res-asignaciones}

\end{table}

\begin{table}[htbp]
\centering
\small
\begin{tabular}{@{}lcccccc@{}}
\toprule
 & \multicolumn{3}{c}{Sin post-proceso} & \multicolumn{3}{c}{B2 (Optuna)}\\
\cmidrule(lr){2-4}\cmidrule(lr){5-7}
Pérdida & Prec. & Exh. & Sin fil.\ (\%) & Prec. & Exh. & Sin fil.\ (\%)\\
\midrule
Baseline (s42) & 0,43 & 0,57 & 27,5 & 0,63 & 0,53 & 9,4\\
Baseline (s43) & 0,36 & 0,58 & 34,4 & 0,61 & 0,50 & 12,4\\
IoU & 0,37 & 0,55 & 29,3 & 0,58 & 0,51 & 12,5\\
clDice & 0,44 & 0,56 & 25,2 & 0,57 & 0,53 & 14,7\\
IoU+clDice & 0,46 & 0,56 & 24,6 & 0,59 & 0,55 & 13,7\\
\midrule
Topograph (s42) & \textbf{0,49} & 0,53 & \textbf{18,0} & 0,57 & 0,53 & 13,5\\
Topograph (s43) & \textbf{0,51} & 0,49 & \textbf{13,4} & \textbf{0,65} & 0,47 & \textbf{7,8}\\
IoU+Topo & \textbf{0,48} & 0,50 & \textbf{15,8} & 0,62 & 0,46 & \textbf{8,3}\\
clDice+Topo & \textbf{0,49} & 0,51 & \textbf{17,4} & 0,59 & 0,50 & 11,1\\
IoU+clDice+Topo & \textbf{0,50} & 0,51 & \textbf{18,2} & 0,57 & 0,49 & 11,2\\
\bottomrule
\end{tabular}
\caption{Fiabilidad de las detecciones en val~$B$. Prec.: aciertos entre
predicciones, $\mathrm{TP}/(\mathrm{TP}+\mathrm{FP})$. Exh.: aciertos
entre filamentos de referencia, $\mathrm{TP}/869$. Sin fil.: porcentaje
de predicciones que no tocan ningún filamento. En negrita, los valores
en los que todas las variantes con Topograph mejoran a todas las
variantes sin él.}
\label{tab:res-fiabilidad}
\end{table}

% ------------------------------------------------------------
\section{Incertidumbre de las comparaciones}
\label{sec:res-bootstrap}
% ------------------------------------------------------------

La Tabla~\ref{tab:res-bootstrap} da la diferencia de PQ de cada
ejecución con la baseline (semilla 42) y su intervalo del 95\,\%,
obtenido por \emph{bootstrap} pareado por días
(Capítulo~\ref{cap:metodologia}). Sin post-proceso, solo Topograph con
la semilla 42 mejora a la baseline de forma significativa; con la
semilla 43 también la mejora ($+0{,}011$), pero su intervalo incluye el
0, e IoU es significativamente peor. Con B2, siete de las nueve
ejecuciones quedan significativamente por debajo.

Estos intervalos deben leerse con cuidado. La baseline con la semilla 43
es significativamente peor que con la 42 sin post-proceso
($-0{,}024$) y casi lo es con B2, a pesar de tratarse de la misma
configuración. El \emph{bootstrap} por días mide cuánto depende el PQ de
los días que forman val~$B$, pero no la variabilidad del propio
entrenamiento. Con una sola ejecución por pérdida, una diferencia
«significativa» puede deberse a que la ejecución de referencia salió
especialmente bien o mal.

\begin{table}[htbp]
\centering
\small
\begin{tabular}{@{}lcccc@{}}
\toprule
 & \multicolumn{2}{c}{Sin post-proceso} & \multicolumn{2}{c}{B2 (Optuna)}\\
\cmidrule(lr){2-3}\cmidrule(lr){4-5}
Pérdida & $\Delta$PQ & IC 95\,\% & $\Delta$PQ & IC 95\,\%\\
\midrule
Baseline (s42) & --- & --- & --- & --- \\
Baseline (s43) & $-0{,}024^*$ & $[-0{,}039,\,-0{,}009]$ & $-0{,}016$ & $[-0{,}035,\,+0{,}005]$ \\
IoU & $-0{,}026^*$ & $[-0{,}044,\,-0{,}010]$ & $-0{,}022^*$ & $[-0{,}037,\,-0{,}007]$ \\
clDice & $+0{,}002$ & $[-0{,}011,\,+0{,}016]$ & $-0{,}019^*$ & $[-0{,}036,\,-0{,}002]$ \\
IoU+clDice & $+0{,}011$ & $[-0{,}004,\,+0{,}024]$ & $-0{,}010$ & $[-0{,}030,\,+0{,}011]$ \\
Topograph (s42) & $+0{,}016^*$ & $[+0{,}001,\,+0{,}032]$ & $-0{,}020^*$ & $[-0{,}035,\,-0{,}001]$ \\
Topograph (s43) & $+0{,}011$ & $[-0{,}008,\,+0{,}029]$ & $-0{,}019^*$ & $[-0{,}035,\,-0{,}003]$ \\
IoU+Topo & $+0{,}004$ & $[-0{,}017,\,+0{,}025]$ & $-0{,}028^*$ & $[-0{,}047,\,-0{,}008]$ \\
clDice+Topo & $+0{,}005$ & $[-0{,}010,\,+0{,}022]$ & $-0{,}025^*$ & $[-0{,}044,\,-0{,}006]$ \\
IoU+clDice+Topo & $+0{,}007$ & $[-0{,}009,\,+0{,}022]$ & $-0{,}037^*$ & $[-0{,}055,\,-0{,}019]$ \\
\bottomrule
\end{tabular}
\caption{Diferencia de PQ en val~$B$ con la baseline (semilla 42) e
intervalo del 95\,\% (\emph{bootstrap} pareado por días, 2.000
remuestreos). $^*$: el intervalo no contiene el 0.}
\label{tab:res-bootstrap}
\end{table}

% ------------------------------------------------------------
\section{Post-proceso por persistencia}
\label{sec:anx-e2}
% ------------------------------------------------------------

E2 y B2+P se aplicaron a ejecuciones reentrenadas con la semilla 42, que son las
ocho pérdidas con la configuración común (\emph{ref.}) y los mejores
intentos del ajuste de hiperparámetros de la baseline, clDice y
Topograph (\emph{mejor}). Las ejecuciones de la Tabla~\ref{tab:res-perdidas}
no se conservaron, y por el no determinismo del entrenamiento las cifras
de los reentrenamientos difieren algo de las suyas. E2 se aplicó a los
once modelos y B2+P también (Tablas~\ref{tab:res-e2}
y~\ref{tab:res-b2p}). Los tipos de error de los cuatro post-procesos se
detallan para seis de ellos (Tabla~\ref{tab:res-e2-errores}).

\paragraph{E2.} Queda por debajo de B2 en PQ en los once modelos, entre
$0{,}03$ y $0{,}06$ sin Topograph y entre $0{,}004$ y $0{,}024$ con
Topograph. En cinco de los seis modelos de la
Tabla~\ref{tab:res-e2-errores}, los filamentos sin predicción quedan al
nivel de la predicción sin post-proceso (B2 los aumenta en entre 48 y
141) y E2 es, de los tres
post-procesos, el que conserva más espina; la excepción es clDice+Topo,
para el que el ajuste elige un umbral alto ($0{,}68$). El ajuste elige
$\delta$ entre $0{,}22$ y $0{,}30$ en nueve de los once modelos, y los
partidos bajan en los seis respecto a la predicción sin post-proceso,
con un coste en fusiones pequeño (los unidos suben como mucho en 7). Los partidos pasan
de 132 a 107 en la baseline de referencia y de 88 a 45 en Topograph de
referencia. En la baseline, donde $\tau$ es casi nulo y el filtro apenas
actúa, la reducción se debe a la fusión; con Topograph, el filtro
elimina además trozos poco persistentes. B2 deja aún menos partidos en
cinco de los seis modelos, pero lo consigue eliminando componentes
pequeñas, entre ellas fragmentos, y con ello también filamentos
enteros.
Para la baseline, el ajuste elige un $\tau$ casi nulo
($0{,}03$--$0{,}06$), y E2 elimina solo un 10--12\,\% de las
predicciones sin filamento; con Topograph elimina entre un 12 y un
38\,\%, frente a un 3--12\,\% en las ejecuciones sin él. Por eso, con la baseline, los errores de presencia
de E2 son muy superiores a los de B2 (421--479 frente a 321--327),
mientras que con Topograph son del mismo orden (321--342 frente a
320--345), aunque con otra composición: E2 pierde menos filamentos y
deja más espurias.

\paragraph{B2+P.} El resultado depende de si el modelo se entrenó con
Topograph (Tabla~\ref{tab:res-b2p}). En los cinco modelos con Topograph,
B2+P iguala o mejora ligeramente el PQ de B2 (hasta $+0{,}012$), y tiene
siempre menos partidos y unidos y mayor precisión. El menor número de
errores de estructura de todo el trabajo, 32, lo da IoU+clDice+Topo con
B2+P. En los seis modelos sin Topograph, en cambio, el PQ queda igual
(baseline mejor) o por debajo, y la precisión baja en todos, porque la
persistencia rescata también espurias muy seguras. Con la baseline de
referencia, por ejemplo, los filamentos sin predicción bajan de 238 a
129, pero las predicciones sin filamento suben de 83 a 233.

El ajuste elige siempre un $\tau$ muy alto ($0{,}88$--$1{,}00$) y un área
mínima mayor que la de B2 (59--283 píxeles frente a 21--42). Con $\tau$
en el máximo del rango, B2+P se reduce prácticamente a B2, así que en
val~$A$ debería igualarlo como mínimo. No ocurre en seis de los once
modelos, cinco de ellos sin Topograph y el mejor intento de Topograph
(por ejemplo, $0{,}343$ frente a $0{,}359$ en la baseline de
referencia). Con tres parámetros y 60 intentos, la búsqueda no siempre
encuentra el óptimo, y parte de la desventaja de B2+P sin Topograph
puede deberse a ello.

\begin{table}[htbp]
\centering
\small
\begin{tabular}{@{}lcccccc@{}}
\toprule
 & \multicolumn{2}{c}{PQ} & \multicolumn{2}{c}{P.\,+\,U.} & \multicolumn{2}{c}{Prec.}\\
\cmidrule(lr){2-3}\cmidrule(lr){4-5}\cmidrule(lr){6-7}
Modelo & B2 & B2+P & B2 & B2+P & B2 & B2+P\\
\midrule
Baseline (ref.)   & \textbf{0,370} & 0,347 & \textbf{59} & 68 & \textbf{0,613} & 0,498\\
Baseline (mejor)  & 0,368 & 0,368 & 49 & \textbf{36} & \textbf{0,632} & 0,600\\
IoU               & \textbf{0,361} & 0,353 & \textbf{47} & 50 & \textbf{0,634} & 0,543\\
clDice (ref.)     & \textbf{0,344} & 0,339 & 55 & \textbf{46} & \textbf{0,561} & 0,516\\
clDice (mejor)    & \textbf{0,361} & 0,354 & 58 & \textbf{49} & \textbf{0,567} & 0,525\\
IoU+clDice        & \textbf{0,363} & 0,347 & 49 & 49 & \textbf{0,596} & 0,516\\
\midrule
Topograph (ref.)  & 0,348 & \textbf{0,360} & 45 & \textbf{37} & 0,616 & \textbf{0,630}\\
Topograph (mejor) & 0,362 & 0,362 & 52 & \textbf{34} & 0,605 & \textbf{0,612}\\
IoU+Topo          & 0,349 & \textbf{0,353} & 57 & \textbf{52} & 0,569 & \textbf{0,575}\\
clDice+Topo       & 0,335 & \textbf{0,341} & 58 & \textbf{35} & 0,574 & \textbf{0,607}\\
IoU+clDice+Topo   & 0,338 & \textbf{0,348} & 58 & \textbf{32} & 0,559 & \textbf{0,611}\\
\bottomrule
\end{tabular}
\caption{B2 frente a B2+P en val~$B$ en los once modelos reentrenados.
P.\,+\,U.: partidos más unidos. Prec.: precisión. En negrita, el mejor
de los dos post-procesos en cada medida cuando difieren.}
\label{tab:res-b2p}
\end{table}

\begin{table}[htbp]
\centering
\small
\begin{tabular}{@{}lccccc@{}}
\toprule
Modelo & Sin post & B2 & E2 & B2+P & $\tau$ (E2)\\
\midrule
Baseline (ref.) & 0,297 & \textbf{0,370} & 0,310 & 0,347 & 0,06\\
Baseline (mejor) & 0,303 & \textbf{0,368} & 0,314 & \textbf{0,368} & 0,03\\
IoU & 0,299 & \textbf{0,361} & 0,313 & 0,353 & 0,42\\
clDice (ref.) & 0,307 & \textbf{0,344} & 0,307 & 0,339 & 0,48\\
clDice (mejor) & 0,320 & \textbf{0,361} & 0,333 & 0,354 & 0,21\\
IoU+clDice & 0,317 & \textbf{0,363} & 0,324 & 0,347 & 0,26\\
\midrule
Topograph (ref.) & 0,316 & 0,348 & 0,337 & \textbf{0,360} & 0,50\\
Topograph (mejor) & 0,330 & \textbf{0,362} & 0,339 & \textbf{0,362} & 0,15\\
IoU+Topo & 0,305 & 0,349 & 0,325 & \textbf{0,353} & 0,50\\
clDice+Topo & 0,299 & 0,335 & 0,332 & \textbf{0,341} & 0,48\\
IoU+clDice+Topo & 0,311 & 0,338 & 0,327 & \textbf{0,348} & 0,35\\
\bottomrule
\end{tabular}
\caption{PQ en val~$B$ de los modelos reentrenados con cuatro
post-procesos: ninguno, B2 (umbral y área), E2 (umbral, persistencia y
fusión por collado) y B2+P (umbral, y área o persistencia), los tres
últimos ajustados en val~$A$. \emph{ref.}: configuración común;
\emph{mejor}: mejor intento del ajuste de hiperparámetros
(Sección~\ref{sec:res-hp}). $\tau$ es la persistencia mínima elegida
para E2. En negrita, el mejor post-proceso de cada modelo.}
\label{tab:res-e2}
\end{table}

\begin{table}[htbp]
\centering
\footnotesize
\setlength{\tabcolsep}{4pt}
\begin{tabular}{@{}llccccccc@{}}
\toprule
Modelo & Post & Partidos & Unidos & Sin pred. & Sin fil. & Prec. & Exh. & Cob.\ espina\\
\midrule
Baseline (ref.) & sin  & 132 & 12 & 97  & 424 & 0,376 & 0,591 & 0,772\\
                & B2   & 47  & 12 & 238 & 83  & 0,613 & 0,532 & 0,699\\
                & E2   & 107 & 13 & 97  & 382 & 0,403 & 0,595 & 0,772\\
                & B2+P & 56  & 12 & 129 & 233 & 0,498 & 0,589 & 0,749\\
\midrule
Baseline (mejor) & sin  & 97  & 15 & 143 & 318 & 0,410 & 0,544 & 0,739\\
                 & B2   & 34  & 15 & 255 & 72  & 0,632 & 0,516 & 0,687\\
                 & E2   & 74  & 17 & 142 & 279 & 0,438 & 0,544 & 0,741\\
                 & B2+P & 21  & 15 & 213 & 116 & 0,600 & 0,539 & 0,688\\
\midrule
Topograph (ref.) & sin  & 88  & 14 & 164 & 199 & 0,458 & 0,533 & 0,715\\
                 & B2   & 31  & 14 & 281 & 64  & 0,616 & 0,480 & 0,669\\
                 & E2   & 45  & 20 & 156 & 165 & 0,519 & 0,541 & 0,733\\
                 & B2+P & 23  & 14 & 260 & 73  & 0,630 & 0,503 & 0,666\\
\midrule
Topograph (mejor) & sin  & 75  & 13 & 189 & 158 & 0,498 & 0,529 & 0,707\\
                  & B2   & 39  & 13 & 237 & 89  & 0,605 & 0,522 & 0,686\\
                  & E2   & 60  & 16 & 191 & 139 & 0,525 & 0,526 & 0,703\\
                  & B2+P & 22  & 12 & 235 & 93  & 0,612 & 0,520 & 0,673\\
\midrule
IoU+Topo & sin  & 88  & 10 & 168 & 233 & 0,431 & 0,520 & 0,702\\
         & B2   & 43  & 14 & 224 & 102 & 0,569 & 0,516 & 0,693\\
         & E2   & 51  & 17 & 165 & 177 & 0,490 & 0,524 & 0,717\\
         & B2+P & 38  & 14 & 227 & 110 & 0,575 & 0,521 & 0,692\\
\midrule
clDice+Topo & sin  & 54  & 25 & 172 & 172 & 0,466 & 0,487 & 0,727\\
            & B2   & 36  & 22 & 234 & 86  & 0,574 & 0,489 & 0,701\\
            & E2   & 33  & 19 & 229 & 107 & 0,557 & 0,490 & 0,686\\
            & B2+P & 13  & 22 & 254 & 78  & 0,607 & 0,480 & 0,670\\
\bottomrule
\end{tabular}
\caption{Tipos de error en val~$B$ de seis modelos reentrenados con los
cuatro post-procesos. Partidos, Unidos, Sin pred.\ y Sin fil.\ como en
la Tabla~\ref{tab:res-asignaciones}; Prec.\ y Exh.\ como en la
Tabla~\ref{tab:res-fiabilidad}. «Cob.\ espina»: fracción de la espina
anotada cubierta por la predicción (Sección~\ref{sec:res-barbas}).}
\label{tab:res-e2-errores}
\end{table}

% ------------------------------------------------------------
\section{Barbas}
\label{sec:anx-barbas}
% ------------------------------------------------------------

De los 869 filamentos de val~$B$, el 41,4\,\% tiene al menos una barba,
con una media de $0{,}70$ barbas por filamento. Su esqueleto mide en
total un 6\,\% de la longitud de las espinas, y el grosor mediano de un
filamento es de $11{,}4$ píxeles a resolución completa (menos de 3 a
$512\times512$), así que el PQ apenas las tiene en cuenta.

Para separar lo que pierde el modelo de lo que se pierde por la
resolución, se reduce cada máscara anotada a $512\times512$, se vuelve a
ampliar y se compara con la original. Se obtiene una cobertura de barbas
de $0{,}997$, una precisión de $0{,}908$ y un $F_1$ de $0{,}950$. La
máscara a $512\times512$ conserva, por tanto, las barbas, y la precisión
de los modelos debe compararse con ese techo de $0{,}91$, no con 1,
porque el reescalado crea pequeñas ramas que la medida cuenta como
inventadas. Queda abierto si la resolución de entrada limita la
detección, ya que a esa escala una barba mide entre 1 y 3 píxeles.

Las medidas se calcularon sobre los reentrenamientos con la
configuración común (Tabla~\ref{tab:res-barbas}). La baseline tiene el
mejor $F_1$ con y sin post-proceso. En las cuatro parejas que solo se
diferencian en clDice, añadirlo baja el $F_1$; sin post-proceso, la
cobertura sube en tres de ellas pero la precisión baja en las cuatro.
Una posible explicación es que el esqueleto de las barbas es solo un
6\,\% del de las espinas, de modo que clDice optimiza sobre todo la
continuidad de la espina. Topograph cubre menos barbas y también menos
espina ($0{,}715$ frente a $0{,}772$), que es el mismo efecto de predecir
menos a cambio de predicciones más fiables. Con B2, la cobertura de la
espina baja más que la de las barbas, lo que indica que el filtro de
área elimina sobre todo filamentos pequeños enteros. Los post-procesos
por persistencia apenas cambian el panorama, ya que E2 deja las medidas
casi como sin post-proceso y B2+P casi como B2. La baseline mantiene el
mejor $F_1$ con los cuatro post-procesos, y añadir clDice lo baja en
todas las parejas y con todos ellos.

\begin{table}[htbp]
\centering
\small
\begin{tabular}{@{}lccccc@{}}
\toprule
Pérdida & Cobertura & Precisión & $F_1$ & Ramas inv. & Cob.\ espina\\
\midrule
Techo a $512\times512$ & 0,997 & 0,908 & 0,950 & 0,075 & 0,999\\
\midrule
\multicolumn{6}{l}{\emph{Sin post-proceso}}\\
Baseline & \textbf{0,563} & \textbf{0,279} & \textbf{0,373} & 0,222 & \textbf{0,772}\\
IoU & 0,536 & 0,268 & 0,357 & 0,226 & 0,741\\
clDice & 0,546 & 0,248 & 0,341 & 0,212 & 0,754\\
IoU+clDice & 0,554 & 0,253 & 0,347 & 0,200 & 0,753\\
Topograph & 0,512 & 0,232 & 0,320 & 0,169 & 0,715\\
IoU+Topo & 0,496 & 0,275 & 0,354 & 0,168 & 0,702\\
clDice+Topo & 0,537 & 0,190 & 0,281 & 0,229 & 0,727\\
IoU+clDice+Topo & 0,529 & 0,240 & 0,331 & 0,185 & 0,728\\
\midrule
\multicolumn{6}{l}{\emph{B2 (Optuna)}}\\
Baseline & 0,528 & \textbf{0,302} & \textbf{0,384} & 0,174 & 0,699\\
IoU & 0,521 & 0,294 & 0,376 & 0,186 & 0,670\\
clDice & \textbf{0,532} & 0,249 & 0,340 & 0,196 & \textbf{0,715}\\
IoU+clDice & 0,516 & 0,266 & 0,351 & 0,174 & 0,704\\
Topograph & 0,504 & 0,229 & 0,315 & 0,155 & 0,669\\
IoU+Topo & 0,508 & 0,275 & 0,357 & 0,171 & 0,693\\
clDice+Topo & 0,497 & 0,202 & 0,287 & 0,217 & 0,701\\
IoU+clDice+Topo & 0,517 & 0,252 & 0,339 & 0,171 & 0,706\\
\midrule
\multicolumn{6}{l}{\emph{E2}}\\
Baseline & \textbf{0,562} & \textbf{0,279} & \textbf{0,373} & 0,223 & \textbf{0,772}\\
IoU & 0,543 & 0,266 & 0,357 & 0,227 & 0,743\\
clDice & 0,551 & 0,244 & 0,338 & 0,200 & 0,756\\
IoU+clDice & 0,515 & 0,253 & 0,339 & 0,199 & 0,744\\
Topograph & 0,543 & 0,239 & 0,332 & 0,184 & 0,733\\
IoU+Topo & 0,524 & 0,270 & 0,356 & 0,189 & 0,717\\
clDice+Topo & 0,474 & 0,208 & 0,289 & 0,189 & 0,686\\
IoU+clDice+Topo & 0,501 & 0,258 & 0,340 & 0,166 & 0,711\\
\midrule
\multicolumn{6}{l}{\emph{B2+P}}\\
Baseline & \textbf{0,563} & \textbf{0,294} & \textbf{0,386} & 0,205 & \textbf{0,749}\\
IoU & 0,529 & 0,270 & 0,357 & 0,193 & 0,711\\
clDice & 0,532 & 0,255 & 0,345 & 0,196 & 0,731\\
IoU+clDice & 0,524 & 0,264 & 0,351 & 0,178 & 0,731\\
Topograph & 0,490 & 0,211 & 0,295 & 0,155 & 0,666\\
IoU+Topo & 0,507 & 0,283 & 0,363 & 0,162 & 0,692\\
clDice+Topo & 0,474 & 0,201 & 0,282 & 0,204 & 0,670\\
IoU+clDice+Topo & 0,512 & 0,257 & 0,342 & 0,163 & 0,680\\
\bottomrule
\end{tabular}
\caption{Medidas de barbas en val~$B$ de los reentrenamientos (semilla
42) con los cuatro post-procesos. «Ramas inv.»: ramas inventadas por filamento. «Cob.\ espina»:
fracción de la espina cubierta. El techo es la propia anotación reducida
a $512\times512$ y vuelta a ampliar.}
\label{tab:res-barbas}
\end{table}

% ------------------------------------------------------------
\section{Ajuste de hiperparámetros}
\label{sec:anx-hp}
% ------------------------------------------------------------

Se ajustaron con Optuna los hiperparámetros de la baseline, Topograph y
clDice (Sección~\ref{sec:ajuste-hp}). Cada búsqueda dispuso de 3 horas.
La baseline completó 15 intentos y se podaron 27; Topograph, 8 y 19, y
clDice, 14 y 18. La Tabla~\ref{tab:res-hp} compara el intento de
referencia (la configuración común) con el mejor intento en val~$A$.

Con B2, el mejor intento de la baseline queda igual que la referencia,
mientras que Topograph sube $0{,}014$ y clDice $0{,}017$, lo que sugiere
que la configuración común se adaptaba mejor a la baseline. En ambas
pérdidas el mejor intento usa una tasa de aprendizaje mayor, y en
Topograph un peso menor que el fijado a priori ($\alpha\approx0{,}015$
frente a $0{,}04$). La selección en $A$ infla las cifras de esa mitad
(clDice pasa de $0{,}384$ en $A$ a $0{,}361$ en $B$). Los intentos de
referencia son, además, un nuevo entrenamiento de cada modelo, y
reproducen el patrón del capítulo, ya que sin post-proceso Topograph supera a
la baseline ($+0{,}019$) y con B2 queda por debajo. En todas las
ejecuciones, la época elegida está entre la 8 y la 20 de las 30
entrenadas; las pérdidas topológicas alcanzan su mejor época algo más
tarde que la baseline (12--20 frente a 8--14), pero siempre lejos de la
última, así que no hay indicios de que falten épocas.

\begin{table}[htbp]
\centering
\small
\begin{tabular}{@{}llccccccc@{}}
\toprule
 & & & & & \multicolumn{2}{c}{Sin post} & \multicolumn{2}{c}{B2}\\
\cmidrule(lr){6-7}\cmidrule(lr){8-9}
Modelo & Intento & lr & $w_d$ & Peso & PQ $A$ & PQ $B$ & PQ $A$ & PQ $B$\\
\midrule
Baseline & ref. & $3{,}0\cdot10^{-4}$ & $10^{-4}$ & 0,50 & 0,296 & 0,297 & 0,359 & \textbf{0,370}\\
 & mejor & $4{,}0\cdot10^{-4}$ & $3\cdot10^{-6}$ & 0,33 & 0,309 & 0,303 & 0,367 & 0,368\\
\midrule
Topograph & ref. & $3{,}0\cdot10^{-4}$ & $10^{-4}$ & 0,040 & 0,307 & 0,316 & 0,350 & 0,348\\
 & mejor & $4{,}0\cdot10^{-4}$ & $3\cdot10^{-6}$ & 0,015 & 0,333 & \textbf{0,330} & 0,368 & 0,362\\
\midrule
clDice & ref. & $3{,}0\cdot10^{-4}$ & $10^{-4}$ & 0,25 & 0,305 & 0,307 & 0,340 & 0,344\\
 & mejor & $7{,}9\cdot10^{-4}$ & $10^{-6}$ & 0,26 & 0,341 & 0,320 & 0,384 & 0,361\\
\bottomrule
\end{tabular}
\caption{Ajuste de hiperparámetros. «Peso» es el de Dice en la baseline,
$\alpha$ en Topograph y el de clDice en clDice; $w_d$ es el
\emph{weight decay}. El mejor intento se elige por el PQ en val~$A$ sin
post-proceso; su PQ en $A$ está inflado por esa selección.}
\label{tab:res-hp}
\end{table}

% ------------------------------------------------------------
\section{Test de la competición}
\label{sec:anx-test}
% ------------------------------------------------------------

Se enviaron al \emph{leaderboard} público los mejores intentos del ajuste
de hiperparámetros. Con B2, la baseline y Topograph obtienen un PQ de
$0{,}27$ y clDice de $0{,}30$, entre $0{,}06$ y $0{,}10$ por debajo de
su PQ en val~$B$. Con B2+P, la baseline y Topograph obtienen $0{,}28$,
y Topograph con E2, $0{,}26$, con una caída respecto a val~$B$ de entre
$0{,}08$ y $0{,}10$ en todas las configuraciones.

Se comprobaron dos posibles causas de la caída. La distribución de
estaciones es equivalente en entrenamiento y test (entre el 14 y el
19\,\% de las imágenes en cada una), y la de años es parecida (un
17\,\% de imágenes de 2018--2022 en el test frente a un 13\,\% en
entrenamiento). Ninguna de las dos explica la diferencia. Que clDice
quede por delante tampoco permite sacar conclusiones, ya que se trata de
un solo envío por modelo y el \emph{leaderboard} público usa solo la
mitad de las imágenes del test.
